\section{The source density, G completion and the distinct sums}
\label{sec:source-colour}

The removed signed measure must correspond to the density that actually
generated the H sample, on its retained component space. A known
nonnegative multiplier may be separated when constructing a ratio, but
must remain in the physical amount removed. This requirement is stronger than
identifying a familiar splitting function. The source includes native
inverse maps, luminosities, colour probabilities, regularized counterevents
and low-scale redistributions. The scalar kernel used to describe a shower
proposal is a separate object. The distinction is essential to applying the
multiplicative--accumulative idea of Ref.~\cite{NasonSalam} to the present
implementation.

\subsection{A common measure and the \texorpdfstring{$S/H$}{S/H} decomposition}
\label{sec:source-measures}

Let $b$ denote a Born configuration and $q$ a labelled real configuration,
including its physical external flavours. Densities below are defined with
respect to positive measures $\dd\mu_B$ and $\dd\mu_R$. A compatible inverse
history $h$ supplies a projection and radiation coordinates,
\begin{equation}
 q=M_h(b_h,r_h),\qquad
 \dd\mu_R=J_h\,\dd\mu_B\,\dd r_h,\qquad J_h>0.
 \label{eq:source-measure}
\end{equation}
Different histories generally have different $b_h$, different luminosities
and different colour bases at the same $q$. The label $h$ will include a
native emitter history and, where necessary, its recoil group. Additional
colour labels will be kept explicit below. A statement that two numerical
terms are equal is meaningful only after their phase-space measures,
flavours and auxiliary sampling laws have also been identified.

A schematic additive matching functional is
\begin{equation}
 \langle O\rangle_{\text{\mcatnlo}}
 =\int\dd\mu_B\,\overline B_D(b)\,\mathcal U_B[O](b)
 +\int\dd\mu_R\,H(q)\,\mathcal U_H[O](q).
 \label{eq:source-functional}
\end{equation}
The $S$ coefficient $\overline B_D$ includes the Born term, the finite
combination of virtual and integrated subtraction contributions, collinear
remnants and the integrated matching redistribution. Writing it informally
as $B+V+\int D$ suppresses the infrared regularization and the ordinary NLO
subtraction terms. In a conventional local notation one would instead
understand a regulated combination of the form
$B+V+I_C+\int(D-C)$, with the appropriate Born projections and factorization
terms. The generated $S$ integrand is not required to be positive. The
standard MC@NLO and FKS references explain the additive and local-subtraction
structures~\cite{FrixioneWebber,FKS,Alwall}; the numerical definitions in this
branch are specified by its generated weights.

For the ordinary, non-$\Delta$ source it is useful to write
\begin{align}
 H_{\rm ord}(q)
 &=\sum_{h,c}P_{h,c}(q)
       \bigl[R_{h,c}(q)-D_{G,h,c}(q)\bigr],
 &\sum_{h,c}R_{h,c}(q)&=R(q),
 \label{eq:source-hordinary}\\
 D_{\rm eff}(q)&\equiv R(q)-H_{\rm ord}(q).
 \label{eq:source-deff}
\end{align}
Here $c$ is a native Born colour flow; $R_{h,c}$ includes the real-emission
partition and its flow/recoil assignments. The symbol $D_{G,h,c}$ denotes
the complete MC contribution of that history, including the G replacements.
Both terms have already been expressed in the common real measure.
$P_{h,c}$ is the source no-emission multiplier. In the ordinary mode it
reduces to the history's flow-independent artificial profile $P_h$
described below. The stochastic
conditional multiplier used by \DeltaMC{} requires a larger state space and
is discussed separately in this paper. Equation~\eqref{eq:source-hordinary}
is a marginal sum: the current sampled implementation retains the Born
diagram $d$ as well as $h,c$ when classifying and processing H events.
It does not substitute the negative part of this marginal for the
negative part of a sampled component.

Here $P_{h,c}$ always means the multiplier actually used by the source
routine. If the native Born counterevent fails generation cuts, the
ordinary routine leaves this multiplier at one and does not evaluate a
matching subtraction for that counterevent. Thus an arbitrary cut process
must not be reconstructed by applying the low-scale profile indiscriminately
to every real history. The source routines' cut decisions are part of
Eq.~\eqref{eq:source-hordinary}; the formula is not a prescription to change
them during postprocessing.

Equation~\eqref{eq:source-deff} is an exact definition of the effective
source subtraction, not an identification with a physical shower emission
kernel. The statement $H=R-K$ is a special matching construction and must
not be inserted into the present rejection formula without checking the
actual source. Likewise, neither $R-H$ nor an individual signed
colour-correlated counterterm is automatically a usable nonnegative
probability denominator.

\subsection{The distinct sector, history, diagram and colour labels}
\label{sec:source-labels}

The implementation contains several different sums. Their roles are most
transparent when the labels are separated as follows.
\begin{center}
\small
\begin{tabular}{@{}p{0.20\linewidth}p{0.74\linewidth}@{}}
\toprule
Label & Meaning \\
\midrule
$\alpha=(i,j)$ & The outer FKS sector or integration owner and its nonnegative
partition $\mathcal S_\alpha(q)$. It also provides the sampled event's stored
owner information. \\
$h=(k,l;\rho)$ & A native splitting history, consisting of the emitted and
emitting FKS legs and, if needed, a native recoil group $\rho$. It supplies
its own underlying Born state, map, kernels and counterevents. \\
$d$ & Born integration configuration and associated amplitude/graph.
Its real and Born projection weights are retained separately. \\
$c$ & A leading Born colour flow in the basis of that native Born process.
It supplies a possible planar colour assignment and associated scales. \\
$\ell$ & A colour or anticolour connection of the emitting Born parton
within the selected flow. A gluon can have two connections. \\
\bottomrule
\end{tabular}
\end{center}

In particular, the $kl$ in
\code{\detokenize{compute_MCsubtraction_kl(k_fks,l_fks,...)}} is a native
emitter history. It is not itself the label of a Born colour flow, and is
not the colour-partner index. The latter is found later by
\code{\detokenize{find_color_connectors}}. The wording ``sum over colour
dipoles'' is therefore insufficient to specify what
\code{\detokenize{MCExplicitKLSum}} does in this branch.

The outer $ij$ sum partitions the real matrix element through the FKS
functions $\mathcal S_{ij}$. The inner $kl$ sum reconstructs the MC matching
contribution using every compatible native projection at the same real
point. A native contribution also contains its own identical-particle
partition, implemented as \code{\detokenize{fks_Hij}} for the MC term; this
must not be confused with the outer \code{\detokenize{fks_Sij}} multiplying
the fully assembled hard remainder. Born colour flows and the connections
inside them are additional sums, not alternative names for these two FKS
labels.

The routine \code{\detokenize{compute_native_NLOPS_weights}} first forms the
real, raw MC and G-completion terms for one native history. The driver can
then call \code{\detokenize{repartition_MC_H}} to replace its outer $H$
records with
\begin{equation}
 \widehat H_\alpha(q)=\mathcal S_\alpha(q)\,
       \sum_{h,c}P_{h,c}(q)
       \bigl[R_{h,c}(q)-D_{G,h,c}(q)\bigr].
 \label{eq:source-outerpartition}
\end{equation}
The nonnegative outer partition multiplies the \emph{whole} completed
native sum. Applying it only to a raw MC kernel would not reconstruct
Eq.~\eqref{eq:source-outerpartition}. Summing the outer partitions gives the
global hard density on the qualified physical flavour space. In sampled
mode, $a=(h,d,c)$ below denotes the retained native component; it is not
the outer integration owner $\alpha$. An S event's original $ij$ owner
is provenance, while its actual shower branch identifies the native
history used in the veto.

\subsection{Sampling a colour flow and summing its contributions}
\label{sec:source-colourflows}

The Born flow selector uses nonnegative leading-flow weights $w_c(b)$ and
probabilities
\begin{equation}
 p_c(b)=\frac{w_c(b)}{\sum_{c'}w_{c'}(b)},\qquad
 \sum_cp_c=1.
 \label{eq:source-colourprob}
\end{equation}
The implementation's \code{\detokenize{get_mbar}} calls this an Odagiri
prescription~\cite{Odagiri}: a flow fraction partitions the full Born and
spin-correlated Born amplitudes among planar flow assignments. Explicitly
summing this chosen flow partition does not make the subsequent shower an
exact quantum colour-density-matrix evolution. The hard matrix elements,
the planar shower colour assignment, and exact soft colour correlations
enter different parts of the construction.

It is important to distinguish the physical fraction $p_c$ from the
probability $q_c$ actually used to draw a flow. In the implemented sampler
$q_c=p_c$, but they play different algebraic roles. At fixed underlying
Born kinematics,
\begin{equation}
 \mathbb E_{c\sim q}\!\left[\frac{X_c}{q_c}\right]
 =\sum_cX_c.
 \label{eq:source-colourmc}
\end{equation}
This is a linear identity. It does not imply an identity after taking a
negative part or an absolute value.

The raw MC Born and spin-correlated Born terms returned by
\code{\detokenize{get_mbar}} already include their flow fraction. The routine
\code{\detokenize{include_born_flow_weight}} therefore supplies only
$1/q_c$ to those raw MC prefactors, whereas the colour-summed real and G
replacement terms receive $p_c/q_c$. Multiplying an additional $p_c$ into
the raw MC term would count its colour fraction twice. For explicit colour
summation the corresponding call uses the physical flow fraction and
sampling probability one. The finite Born/virtual contributions are
included once, rather than once per flow.

Even the statement ``the flow probability'' can hide a scale dependence.
The source selects the native flow at its native Born point and coupling
scale. The raw MC Born fraction can be evaluated at the real-state coupling
scale. In processes with mixed Born coupling orders the two fractions need
not be interchangeable. A sampled probability must be retained with the
draw; recomputing its denominator at a different projection or scale is not
an equivalent operation. The summed-colour routine deliberately includes
leading flows whose selected-point fraction is zero, since another
contribution's coupling convention can give a different flow fraction.

Each shower event retains one actual colour assignment. Setting
\code{\detokenize{MCSubtractionAtFixedFlow}} to false does not mean that a
shower event contains all colours simultaneously. The source evaluates and
combines the subtraction weights over native Born flows, while retaining
the sampled \emph{outer owner} flow and its event colour/scale information.
Each inner history uses its own Born colour basis. Reusing the outer flow's
integer label inside another native Born process would be invalid unless
an explicit map between those bases had first been supplied.

Within a fixed flow the connection loop must also preserve multiplicity.
For a gluon connected by colour and anticolour to the same partner, the
code restores two connections even though they share one partner index.
Dropping one because the indices coincide would change the kernel.

\subsection{The G function is an additive completion}
\label{sec:source-gfunction}

A scalar shower approximation has a useful collinear limit, but its soft
colour structure, spin dependence and angular support need not coincide
point by point with the limiting expressions required by a local FKS
subtraction. The G construction interpolates between these expressions. In
the present code it is not equivalent to multiplying the complete MC
counterterm by a positive scalar damping factor.

For one native history and flow, its algebraic structure can be represented
as
\begin{align}
 D_{G,h,c}={}&\sum_{\ell}d_{h,c,\ell}\,g_{s,h}
       \bigl(K^{\rm sc}_{h,c,\ell}
             +g_{\phi,h}K^{\rm spin}_{h,c,\ell}\bigr)
       \nonumber\\
 &+(1-\widetilde g_{s,h,c})\,T_{s,h,c}
       \nonumber\\
 &+(1-\widetilde g_{s,h,c})(1-g_{c,h})
       \bigl(T_{c,h,c}-T_{sc,h,c}\bigr),
 \label{eq:source-gcompletion}\\
 \widetilde g_{s,h,c}
 ={}&1-(1-g_{s,h})\,\overline d_{h,c},
 &\overline d_{h,c}
 ={}&\frac{1}{n_{h,c}}\sum_{\ell=1}^{n_{h,c}}d_{h,c,\ell}.
 \label{eq:source-effectiveg}
\end{align}
Here $d_{h,c,\ell}$ is the connection's starting-scale damping, and
$K^{\rm sc}$ and $K^{\rm spin}$ denote the scalar and azimuthally
spin-correlated raw MC terms. The symbols $T_s,T_c,T_{sc}$ stand for the
soft, collinear and soft-collinear replacement templates, including their
own measures, luminosities and prefactors. They are not the $S$ event class
and they are not simple functions of the real momenta alone. Equation
\eqref{eq:source-gcompletion} displays the source algebra, rather than
replacing the native template definitions by a universal closed formula.

There are two different soft factors in these equations. The raw kernels
use $g_s$. Only after their connection loop does
\code{\detokenize{compute_MCsubtraction_kl}} return the effective
$\widetilde g_s$ used by the additive templates. The connection average in
Eq.~\eqref{eq:source-effectiveg} includes connections outside angular
shower support: the soft completion must still cover wide-angle points
where an individual raw shower term vanishes. The raw connection sum and
the average entering $\widetilde g_s$ must not be confused.

In \code{\detokenize{compute_gfun}},
\begin{equation}
 g_s=G(1-\xi; a_s,b_s,2),\qquad
 g_c=G(y;a_s,-0.1,1),\qquad
 g_\phi=1-G(y;-a_\phi,b_\phi,\delta)
 \label{eq:source-garguments}
\end{equation}
for the enabled azimuthal option. In this notation $\xi$ and $y$ are the
native FKS variables, not laboratory observables. The routine comment
identifies $g_s\to0$ in the soft limit and $g_c\to0$ in the collinear limit.
The standalone evaluator initializes
$(a_s,b_s,a_\phi,b_\phi)=(1,-0.1,-1,-0.1)$, matching the source driver's
input defaults.

For completeness, the function implemented by
\code{\detokenize{gfunction}} has the following shape. For $a\leq0$ it is one.
For $a>0$, define
\begin{equation}
 w_{\min}=\begin{cases}0,&b<0,\\\max(0,1-\delta),&b\geq0,\end{cases}
 \quad w_g=\min\!\bigl(1-(1-w_{\min})|b|,1-10^{-5}\bigr),
 \quad t=\frac{|w|-w_g}{1-w_g}.
 \label{eq:source-gshape-vars}
\end{equation}
The code returns one below the transition, uses
$(1-t)^{2a}/[t^{2a}+(1-t)^{2a}]$ in its transition interval, and sets the
function to zero at $|w|\geq0.99$. The latter is an explicit endpoint
cutoff in the numerical implementation. The word ``smooth'' refers to the
interpolation function; it must not be taken as an assertion that this
additional endpoint cutoff has been removed. These details matter when
attempting to reproduce the counterterm numerically rather than only its
singular limits.

The implementation applies the scalar and spin G factors in
\code{\detokenize{xmcsubt_connection}}. It then adds the templates once in
\code{\detokenize{compute_native_NLOPS_weights}}, with coefficients
$(1-\widetilde g_s)P$ for the soft term and
$(1-g_c)(1-\widetilde g_s)P$ for the collinear and soft-collinear terms, with
their appropriate relative signs. Native counter/real measure ratios are
retained. One cannot reconstruct this density by multiplying a raw
splitting function by $G$, or by evaluating a soft template at the wrong
Born projection.

This also explains why a positive shower denominator does not enforce an
upper bound on the negative hard remainder. A G completion can survive
where a raw scalar history is absent or small. Exact soft colour and spin
terms need not have the same sign or size as the configured scalar shower.
There is no general pointwise inequality $D_G\leq K$, and no guarantee that
$\pos{D_G-R}\leq K$. A required veto ratio above one can therefore be a
property of these different densities and their supports, even after
numerical and bookkeeping errors have been corrected.

\subsection{The artificial ordinary-source no-emission factor}
\label{sec:source-bogus}

For the ordinary non-$\Delta$ source,
\code{\detokenize{bogus_probne_fun}} has mode two enabled and returns
\begin{equation}
 P_{\rm bog}(q_h)=
 \begin{cases}
  0,&q_h\leq0.5\GeV,\\[1mm]
  \displaystyle\frac{u_h^4}{u_h^4+(1-u_h)^4},
    &0.5\GeV<q_h<10\GeV,\\[2mm]
  1,&q_h\geq10\GeV,
 \end{cases}
 \qquad u_h=\frac{q_h-0.5\GeV}{9.5\GeV}.
 \label{eq:source-bogus}
\end{equation}
The argument is that history's \code{\detokenize{get_qMC}} scale. It is a
relative splitting transverse scale and is not in general the measured
jet $\pt$, the total-system recoil, or the same value in all inverse
histories at one $q$. In particular, coloured final states can have
competing ISR and FSR histories with very different values of this scale.

The name ``bogus Sudakov'' distinguishes this fixed profile from a physical
Sudakov no-emission probability. It is not obtained by exponentiating an
integral of the configured shower kernel. Its transition scales are fixed
in GeV, and it is not a factor that approaches one as
$1+O(\as)$ at every fixed point when the coupling is sent to zero.

For a single native component, its paired redistribution has the simple
algebra
\begin{equation}
 H=P(R-D_G),\qquad
 D_{\rm eff}=P D_G+(1-P)R,
 \qquad H+D_{\rm eff}=R.
 \label{eq:source-bogussplit}
\end{equation}
The ordinary $S$ stream contains both $P D_G$ and the transferred
$(1-P)R$, in addition to the other NLO terms. The same $P$ is applied to the
raw MC, real and G replacement pieces. Multiplying only the negative hard
term, or only the post-shower veto denominator, does not implement
Eq.~\eqref{eq:source-bogussplit}.

For example, relative to a source with $P=1$, define the transferred native
density $T_h=(1-P_h)(R_h-D_{G,h})$. Before showering its change of an
observable has the form
\begin{equation}
 \delta_P\langle O\rangle
 =\sum_h\int\dd\mu_R\,T_h(q)
       \bigl[O(\pi_hq)-O(q)\bigr],
 \label{eq:source-bogusobservable}
\end{equation}
where $\pi_h$ is the appropriate native Born projection, with the colour
and other labels understood. This vanishes for $O=1$, but is not an exact
identity for every differential distribution. It can produce low-scale
power effects for suitable infrared-safe observables. In the present
comparisons it is already part of the common ordinary source, so it must
not be invoked as an input difference introduced only by \macnlops.

Different histories have different $P_h$. Thus the hard negative density
is
\begin{equation}
 H_-(q)=\pos{\sum_{h,c}P_{h,c}(q)
                       [D_{G,h,c}(q)-R_{h,c}(q)]},
 \label{eq:source-negativeglobal}
\end{equation}
with the sum taken before the positive part. Pulling out one event's
outer $P_a$, applying a positive-part operation to individual histories,
or multiplying the physical shower denominator by $P_a$ would in general
change the prescription.

Physical \DeltaMC{} uses \code{\detokenize{compute_delta}} in place of
\code{\detokenize{bogus_probne_fun}}; it does not use a product of the two.
The legacy option \code{\detokenize{macnlops_resolved_source=True}} is also
a different source construction. It replaces the ordinary G/P source by
an explicitly gated scalar construction. Restoring an artificial profile
to that altered source would not by itself make it the ordinary MC@NLO
source. The completed campaign uses
\code{\detokenize{macnlops_resolved_source=False}} throughout.

\subsection{Source partition and sign classification before unweighting}
\label{sec:source-sign}

Linear summation and nonlinear sign classification do not commute:
\begin{equation}
 \sum_a\pos{-H_a(q)}\ne
 \pos{-\sum_aH_a(q)}.
 \label{eq:source-signnonlinear}
\end{equation}
For example, contributions $+3$ and $-2$ give a positive total $+1$.
Sampling them separately produces negative events; unweighting their
combined density does not. A postprocessor that receives the separately
sampled $-2$ event cannot treat it as a draw from a globally negative
density that is actually zero. Improving a numerical estimate of the
global denominator does not repair that source-partition mismatch.

The global-source realization therefore forms the complete native $H$ sum and
coalesces compatible physical-flavour contributions before signed $H$
unweighting. This includes three distinct requirements. Every compatible
native history must be available. Every native flow contribution must be
included with its correct fraction. Contributions being combined must
have the same physical event and compatible stored shower ownership,
including colours, scales and, for Delta, conditional context. A sum over
integer FKS labels alone does not establish any of these properties.

The generated registry \code{\detokenize{mc_histories.json}} and the Born
support metadata describe the native mappings and physical flavour rows.
The static completeness certificates and the dynamic H-group test are
both needed. The dynamic test in \code{\detokenize{add_write_info}} checks
compatible H records at the same momenta, with the same colour and scale
ownership and conditional context, before certifying the selected flavour
group. Physical PDG rows are compared by their content rather than by an
ordinal row number which can differ between generated subprocesses.
In the sampled realization the corresponding group is complete only
within one retained native component. Distinct component draws cannot
be merged merely because their momenta coincide. This certifies its
conditional sign, without claiming the sign of the global sum.

The hard annotation then records
\code{\detokenize{partition="global_sum"}} together with effective KL/flow
flags, flavour coalescing and a positive group size. A positive group size
by itself is not a proof of completeness. The importer
\code{\detokenize{hard_source_contract}} checks every input event and
requires a conditional Delta replay packet for a globally summed
Delta $H$ event. Sampled inputs now require their separate qualified
component packet and evaluator (Sec.~\ref{sec:delta-sampled}). The
importer does not infer a global density from the mere presence
of an \code{\detokenize{#aMCatNLO}} sector tag.

In evaluating Eq.~\eqref{eq:source-outerpartition},
\code{\detokenize{repartition_MC_H}} divides each \emph{complete native H
weight} by its positive native real measure and inserts the outer real
measure and $\mathcal S_a$. This conversion retains the counter/real
measure ratios inside the G terms; adding a further Jacobian to a raw
kernel would double count a transformation. Native recoil groups are
also summed, with their partitions evaluated at the appropriate real or
Born point. If a required map is numerically invalid, the generation
routine rejects the complete folded integration point rather than
keeping a partial native sum. Postprocessing a previously generated
source instead stops on an unresolved input; it must not silently omit
that event.

\subsection{Effective FKS settings and the supported combinations}
\label{sec:source-settings}

Three controls have distinct purposes:
\begin{itemize}
 \item \code{\detokenize{MCExplicitKLSum}} enables the complete native
       $kl$ reconstruction of the hard contribution before its outer
       FKS partition is applied.
 \item \code{\detokenize{MCSubtractionAtFixedFlow}} chooses sampled native
       Born-flow subtraction when true, and explicit Born-flow summation
       when false. This is independent of KL summation.
 \item \code{\detokenize{FKSExplicitSum}} sums outer FKS configurations in
       the Born-like evaluation, using a common Born chart where required,
       instead of sampling an allowed configuration group. The Born
       contribution is included once and compatible S terms are combined
       before absolute values are used for unweighting.
\end{itemize}
The last choice can change negative-$S$ fractions and efficiency. It is
not the switch that certifies a complete native hard density. Separate
initial/final-state integration channels remain separate where the
generator treats them separately.

The current postprocessor has the following compatibility
table for the \emph{effective source construction}. ``Either'' in the last
column means true or false, subject to the phase-space and completeness
conditions described here.
\begin{center}
\small
\begin{tabular}{@{}cccp{0.49\linewidth}@{}}
\toprule
KL sum & Fixed flow & Outer FKS sum & \macnlops{} contract \\
\midrule
false & true  & either & Sampled component evaluator; complete
                         diagram/flow/sector provenance required. \\
true  & false & either & Supported when native/flavour completeness and
                         the annotation contract hold. \\
true  & true  & either & Sampled colour source; its conditional negative
                         measure is not implemented by this importer. \\
false & false & either & Native-sector source; a global density cannot
                         replace its sampled negative measure. \\
\bottomrule
\end{tabular}
\end{center}

The shipped FKS card has KL sum false, fixed flow true and outer FKS
sum true. Thus the first two choices already select the sampled
component organization; set the third explicitly to false to sample
the outer groups as well. That setting samples allowed contribution
groups, which may contain associated nonsingular sectors, rather than
guaranteeing a single $ij$ term in every function call. None of these
general generator defaults was changed by the MAcNLOPS implementation.

During the campaign, \code{\detokenize{driver_mintMC.f}} forced
\begin{equation}
 \text{\code{\detokenize{MCExplicitKLSum}}}=\mathrm{true},\qquad
 \text{\code{\detokenize{MCSubtractionAtFixedFlow}}}=\mathrm{false}
 \label{eq:source-forcedflags}
\end{equation}
for \macnlops{} generation. The new driver honours the supported card
pairs above and rejects the other combinations. It does not force
\code{\detokenize{FKSExplicitSum}}. Effective runtime values and the recorded
annotation, rather than a template's comments or module initialization
values, therefore determine compatibility. The campaign explicitly sets
$(\mathrm{KL},\mathrm{fixed\ flow},\mathrm{outer\ sum})=(\mathrm{true},
\mathrm{false},\mathrm{true})$. The outer-sum option requires the common
Born/event-projection construction; it is not an option for UNLOPS in this
implementation. The supported \macnlops{} workflow further requires an
unmerged QCD source, \code{\detokenize{ickkw=0}}, and the declared
\pythia{} setup. Generator support for other modes elsewhere in the branch
is not a validation of \macnlops{} for those modes.

A single flow reduces the local work, once its conditional measure is
specified. The current component evaluator supplies that measure by
keeping the source diagram and flow and identifying the actual shower
history. It stays within the exact source Born provider and transports
diagram and flow indices under certified identical-leg permutations.
An outer event's colour integer alone would not establish this map:
different native Born contexts can use different bases. Retaining the
joint component information is what makes the simpler local evaluation
valid. The colour/anticolour connection sum inside a gluon flow remains.

There are two legitimate constructions: reconstruct and coalesce the global
source before sign classification, or formulate matching on the larger
joint space of momenta, colour and source-owner variables using their
actual conditional laws. Simply exchanging the two boolean settings, or
applying a globally summed ratio to already sampled negative components,
implements neither construction. Section~\ref{sec:delta-sampled}
describes the new qualification and implementation on the joint space.
Anonymous sampled events still lack the information required by that
conditional evaluator, even when only one colour flow is possible.

The principal source anchors lie in
\code{\detokenize{Template/NLO/SubProcesses/}}:
\begin{itemize}
 \item \code{\detokenize{mc_matching_runtime.f}}: sampled component
       capture or complete native H repartition;
 \item \code{\detokenize{montecarlocounter.f}}: raw kernels, connections,
       G factors and source no-emission multipliers;
 \item \code{\detokenize{fks_singular.f}}: real/MC/G weight records and
       flow prefactors;
 \item \code{\detokenize{driver_mintMC.f}}: effective-generation settings;
 \item \code{\detokenize{weight_lines.f}}: component contexts,
       contribution grouping and common Born-state checks;
 \item \code{\detokenize{add_write_info.f}}: physical flavour
       certificates and hard/component packet writing.
\end{itemize}
